% Compatibility entry point: the merged conference-style review lives in related_work.tex.
Machine learning has increasingly been investigated as a non-imaging approach for predicting gallstone disease from routinely available clinical, biochemical, and body composition measurements. Existing studies using the Gallstone-1 dataset have followed several methodological directions, beginning with conventional machine learning classifiers and subsequently extending toward feature selection, metaheuristic optimization, ensemble learning, calibration, and explainable artificial intelligence. Although later studies report improvements in predictive performance, differences in feature selection and evaluation procedures make direct comparison between reported results difficult.

\subsection{Conventional Machine Learning Approaches}

The Gallstone-1 dataset was introduced by Esen et al.~\cite{esen2024early} to investigate whether gallstone disease could be predicted without relying directly on medical imaging. The dataset contains 319 participants, comprising 161 patients with gallstone disease and 158 healthy controls, and includes 38 demographic, laboratory, anthropometric, and bioimpedance variables. The study evaluated a range of conventional machine learning algorithms, including Logistic Regression, Support Vector Machine, K-Nearest Neighbours, Decision Tree, Random Forest, Gradient Boosting, AdaBoost, Naive Bayes, and Multi-Layer Perceptron.

Among the evaluated classifiers, Gradient Boosting produced the highest reported accuracy of 85.42\%~\cite{esen2024early}. Feature analysis also identified Vitamin D, C-reactive protein, total body water, and lean mass as important predictors. These findings established the initial machine learning benchmark for the dataset and showed that gallstone status can be predicted from structured non-imaging measurements. More importantly, the study demonstrated that nonlinear ensemble methods can capture useful relationships among clinical and body composition variables that may not be represented adequately by simpler decision boundaries.

The original study therefore provides two important observations for subsequent work. First, tree-based ensemble models appear suitable for this type of heterogeneous tabular clinical data. Second, only a subset of the available variables appears to provide strong discriminatory information, motivating later work on feature selection and dimensionality reduction.

\subsection{Feature Selection and Metaheuristic Optimization}

Later studies investigated whether optimization and feature selection could improve the original machine learning models or reduce the number of variables required for prediction. Sarker et al.~\cite{sarker2025gallstone} combined Random Forest with the Sand Cat Swarm Optimization algorithm. Four configurations were investigated: Random Forest without cross-validation, Random Forest with cross-validation, optimized Random Forest without cross-validation, and optimized Random Forest with cross-validation.

Using all 38 features, Random Forest without cross-validation achieved an accuracy of 81.25\%, an F1-score of 79.07\%, a precision of 85.00\%, and a recall of 73.91\%~\cite{sarker2025gallstone}. Under 10-fold cross-validation, the same classifier achieved a mean accuracy of 78.42\%. Sand Cat Swarm Optimization reduced the feature set from 38 to 13 variables while maintaining similar predictive performance. The optimized model achieved 79.17\% accuracy without cross-validation and a mean accuracy of 78.32\% under cross-validation.

The contribution of this approach is therefore more evident in dimensionality reduction than in predictive improvement. The optimized feature subset contained approximately one third of the original variables while producing only a small reduction in classification accuracy. SHAP analysis identified CRP, Vitamin D, and aspartate aminotransferase as the most influential predictors, while DiCE was used to provide counterfactual explanations for individual predictions~\cite{sarker2025gallstone}. This study suggests that reducing the number of clinical measurements may be possible without substantially reducing predictive performance.

A related optimization approach was proposed by Shrestha et al.~\cite{shrestha2026metaheuristic}, who applied Rotation Forest and subsequently combined it with the Bald Eagle Search algorithm. The baseline Rotation Forest model using the complete feature set achieved an accuracy of 78.00\%, an F1-score of 77.35\%, precision of 79.33\%, recall of 75.79\%, and an AUC of approximately 0.867. Bald Eagle Search reduced the feature set to 17 variables, but the resulting model achieved a lower accuracy of 75.79\% and an AUC of approximately 0.860~\cite{shrestha2026metaheuristic}.

This result is important because it demonstrates that feature reduction does not necessarily improve predictive accuracy. Although the optimized model used fewer variables, its predictive performance was slightly lower than that of the baseline Rotation Forest. The training results also reached approximately 100\% across folds while test performance remained considerably lower. This difference indicates that complex ensemble methods can fit the small training dataset extremely well without achieving equivalent generalization on unseen observations. SHAP and LIME analyses nevertheless produced findings consistent with previous work, identifying CRP, Vitamin D, obesity-related measurements, haemoglobin, and body mass as influential variables~\cite{shrestha2026metaheuristic}.

Das et al.~\cite{das2026sloa} investigated a similar feature optimization problem using CatBoost combined with the Sea Lion Optimization Algorithm. The baseline CatBoost classifier was evaluated using five-fold cross-validation and achieved a mean accuracy of 79.58\%, an F1-score of 79.01\%, precision of 80.91\%, and recall of 77.36\% using all 38 features. After optimization, the number of features was reduced to 19. The optimized model achieved a mean accuracy of 80.42\%, an F1-score of approximately 79.94\%, precision of 81.97\%, and recall of 77.76\%~\cite{das2026sloa}.

The improvement produced by optimization was therefore relatively small, with mean accuracy increasing by less than one percentage point. However, the feature space was reduced by half. Individual fold performance also varied substantially, with optimized test accuracy ranging from approximately 71.88\% to 87.50\%. This variation further illustrates the sensitivity of performance estimates when the available sample is small. SHAP, LIME, and DiCE were used for model interpretation, with CRP and Vitamin D again identified as dominant predictors~\cite{das2026sloa}.

Taken together, these optimization studies indicate that metaheuristic feature selection can substantially reduce the number of variables required by the models. However, the resulting improvement in predictive performance is inconsistent. Sand Cat Swarm Optimization maintained similar Random Forest performance using fewer predictors, Bald Eagle Search reduced dimensionality but slightly decreased Rotation Forest performance, and Sea Lion Optimization provided only a modest improvement to CatBoost. The primary benefit of these techniques therefore appears to be model simplification rather than a consistent increase in classification accuracy.

\subsection{Ensemble Feature Selection and Model Optimization}

Rather than relying on a single optimization algorithm, Tamer et al.~\cite{tamer2026novel} proposed an ensemble feature selection framework that combines multiple statistical and metaheuristic approaches. The statistical methods included minimum redundancy maximum relevance, Chi-square testing, analysis of variance, and the Kruskal-Wallis test. These were combined with Grey Wolf Optimization, Whale Optimization Algorithm, Harris Hawks Optimization, and Particle Swarm Optimization.

Features selected consistently by at least seven of the eight methods were retained, reducing the original 38 variables to 11~\cite{tamer2026novel}. The selected feature set was subsequently evaluated using Logistic Regression, Naive Bayes, Support Vector Machine, Decision Tree, K-Nearest Neighbour, Gradient Boosting, AdaBoost, Random Forest, and Multi-Layer Perceptron.

Random Forest achieved the highest reported accuracy of 90.62\% using only the 11 selected predictors~\cite{tamer2026novel}. This result is higher than the accuracy reported by the original Gradient Boosting model and by the individual metaheuristic optimization studies. The findings suggest that agreement between multiple feature selection approaches may provide a more stable method of identifying informative predictors than relying on a single optimization algorithm. The authors also reported that clinically meaningful variables, including CRP, haemoglobin, Vitamin D, and hyperlipidaemia, remained within the selected subset.

The study provides evidence that redundant variables may negatively affect prediction when the number of observations is limited. However, the reported improvement should be interpreted together with the feature selection and evaluation procedure. Because the dataset contains only 319 observations, feature selection conducted using information outside the training partition could produce optimistic estimates. Consequently, the effectiveness of the selected feature subset would be strengthened by evaluating it under repeated or nested validation procedures.

Rahman~\cite{rahman2026hbb} approached the optimization problem from a different direction by focusing on hyperparameter optimization rather than only feature subset selection. The proposed HBB-Opt framework combines the global search capability of the Bees Algorithm with Bayesian optimization and applies the resulting strategy to XGBoost.

The optimized XGBoost model achieved an accuracy of 87.50\%, an F1-score of 87.49\%, and an ROC-AUC of 0.9258~\cite{rahman2026hbb}. Sensitivity reached 90.62\%, while the negative predictive value was 90.00\%. The model also achieved a Brier score of 0.098, providing information about probability calibration in addition to discrimination performance. SHAP analysis identified CRP, Vitamin D, obesity percentage, visceral fat area, and the extracellular fluid to total body water ratio as the most influential predictors.

Compared with feature selection approaches, this work demonstrates that improvements can also be obtained by optimizing model parameters and considering calibration. This is particularly relevant in clinical prediction, where the usefulness of a model depends not only on whether observations are classified correctly but also on whether predicted probabilities represent meaningful estimates of disease risk.

\subsection{Ensemble Learning, Calibration, and Explainability}

A further methodological direction involves combining predictions from multiple classifiers. Indirani et al.~\cite{indirani2026ensemble} evaluated ensemble architectures constructed from Random Forest, XGBoost, Logistic Regression, Gradient Boosting, Extra Trees, LightGBM, and Support Vector Machine. Their experiments considered stacking, soft voting, SHAP-based feature selection, probability calibration, and decision threshold optimization.

A stacking model using Random Forest and XGBoost as base learners and Logistic Regression as the meta-learner achieved an accuracy of 85.94\% and an ROC-AUC of 0.9307~\cite{indirani2026ensemble}. The high AUC indicates strong class discrimination and represents one of the highest reported AUC values among studies using this dataset.

The authors also evaluated a broader soft-voting ensemble. At the default classification threshold, the model achieved an accuracy of 79.69\%. After adjusting the classification threshold based on validation performance, accuracy increased to 87.50\%, while recall for gallstone-positive cases reached approximately 94\%~\cite{indirani2026ensemble}. A SHAP-based feature selection variant produced an ROC-AUC of approximately 0.9238 and reached the same adjusted accuracy of 87.50\% following threshold optimization.

These results highlight an aspect that is less visible in studies focusing only on model accuracy. In a clinical classification problem, the decision threshold can substantially change the balance between false-positive and false-negative predictions. A model with high overall accuracy may still be unsuitable for screening if it fails to identify a substantial proportion of patients with the disease. The increase in positive-class recall reported by Indirani et al.~\cite{indirani2026ensemble} therefore demonstrates the importance of considering sensitivity and threshold calibration alongside conventional accuracy metrics.

\subsection{Synthesis of Previous Findings}

Several consistent patterns can be identified across the previous studies. First, tree-based ensemble models have generally produced strong results on the Gallstone-1 dataset. The original study reported 85.42\% accuracy using Gradient Boosting~\cite{esen2024early}, while subsequent work investigated Random Forest~\cite{sarker2025gallstone}, Rotation Forest~\cite{shrestha2026metaheuristic}, CatBoost~\cite{das2026sloa}, XGBoost~\cite{rahman2026hbb}, and heterogeneous ensemble models~\cite{indirani2026ensemble}. This pattern suggests that nonlinear interactions between clinical and biochemical variables are relevant to gallstone prediction.

Second, feature selection is a recurring component of later studies. Sand Cat Swarm Optimization reduced the feature set to 13 variables~\cite{sarker2025gallstone}, Bald Eagle Search retained 17~\cite{shrestha2026metaheuristic}, Sea Lion Optimization retained 19~\cite{das2026sloa}, and the ensemble feature selection method of Tamer et al.~\cite{tamer2026novel} reduced the dataset to 11 predictors. However, the relationship between dimensionality reduction and predictive performance is not consistent. Some approaches improved accuracy, whereas others produced comparable or slightly lower performance. This indicates that the value of feature selection should be assessed not only by the resulting accuracy but also by model complexity, stability, and generalization.

Third, a number of predictors appear repeatedly across independent studies. CRP and Vitamin D are consistently identified as important variables by the original analysis and by subsequent SHAP and feature selection studies~\cite{esen2024early,sarker2025gallstone,shrestha2026metaheuristic,das2026sloa,rahman2026hbb}. Body composition and obesity-related measurements are also frequently retained. The repeated identification of these variables across different modelling strategies provides stronger evidence of their predictive relevance than the importance ranking obtained from any single model.

Despite increasingly complex modelling approaches, the existing literature presents several methodological limitations. The most important limitation is that the majority of studies evaluate their methods on the same 319 observations. Consequently, higher performance reported by a newer model does not necessarily indicate better generalization to an independent clinical population. Performance estimates are also obtained using different validation procedures, feature subsets, optimization methods, and decision thresholds. Therefore, values such as the 90.62\% accuracy reported by the ensemble feature selection approach~\cite{tamer2026novel}, the 87.50\% accuracy reported by HBB-Opt~\cite{rahman2026hbb}, and the 87.50\% threshold-adjusted accuracy reported by the ensemble framework~\cite{indirani2026ensemble} should not be interpreted as direct model rankings.

Another concern is the difference between training and test performance reported in several studies. Random Forest, Rotation Forest, and CatBoost configurations reached or approached perfect training performance while achieving substantially lower validation or test results~\cite{sarker2025gallstone,shrestha2026metaheuristic,das2026sloa}. This behaviour suggests that overfitting remains an important issue for the Gallstone-1 dataset. The relatively small number of observations compared with the available predictors further increases this risk.

Overall, previous research demonstrates that clinical, biochemical, and bioimpedance measurements contain sufficient information to support machine learning based gallstone prediction. Conventional boosting methods established strong initial performance, while subsequent research has explored feature reduction, metaheuristic optimization, ensemble learning, explainability, calibration, and threshold adjustment. However, the reported literature does not demonstrate that increasing algorithmic complexity consistently produces better generalization. In several cases, simpler models perform comparably to optimized alternatives, while feature selection mainly reduces dimensionality.

These findings motivate an evaluation framework in which multiple model families are compared under the same data partitions, preprocessing procedures, feature selection conditions, and evaluation metrics. Such a design allows differences in performance to be attributed more reliably to the models themselves rather than to differences in experimental settings. In addition to accuracy and ROC-AUC, sensitivity, false-negative behaviour, F1-score, and model stability should be considered because these measures are directly relevant to medical classification. Most importantly, results obtained from the Gallstone-1 dataset should be interpreted as evidence of predictive feasibility until the models are validated using larger and independent patient cohorts.
```